\chapter{Phân tích và thiết kế hệ thống}
\label{chap:phan-tich-thiet-ke}

\section{Tác nhân và ca sử dụng}

Tác nhân chính là người quản trị hoặc sinh viên phụ trách phòng lab. Các ca sử dụng gồm quản lý thiết bị, kiểm tra kết nối, mở/đóng tab session, đồng bộ cấu hình, chỉnh desired state, xem trước và đẩy cấu hình, xem backup, duyệt database và tùy chỉnh ứng dụng.

\section{Yêu cầu chức năng}

\begin{longtable}{p{0.12\textwidth}p{0.55\textwidth}p{0.22\textwidth}}
\caption{Yêu cầu chức năng và trạng thái hiện tại}\label{tab:functional-requirements}\\
\toprule
Mã & Yêu cầu & Trạng thái \\
\midrule
\endfirsthead
\toprule
Mã & Yêu cầu & Trạng thái \\
\midrule
\endhead
FR-01 & CRUD, tìm kiếm và nhập danh sách thiết bị từ JSON/XLSX & \statusdone \\
FR-02 & Ping, kết nối, giữ session theo tab và backup running-config & \statusdone \\
FR-03 & Đồng bộ interface và OSPF từ dữ liệu thu thập & \statuspartial \\
FR-04 & CRUD Interface L3/WAN/Tunnel/QoS & \statusdone \\
FR-05 & View \& Push Interface & \statusplanned \\
FR-06 & CRUD và View \& Push DHCP & \statusdone \\
FR-07 & CRUD và View \& Push Static/OSPF/EIGRP & \statuspartial \\
FR-08 & CRUD ACL nhiều loại và binding interface & \statusdone \\
FR-09 & View \& Push ACL & \statusplanned \\
FR-10 & CRUD và View \& Push NAT/PAT & \statusdone \\
FR-11 & Database browser, external tools, theme và status settings & \statusdone \\
FR-12 & VLAN, BGP, VRF và topology automation & \statusplanned \\
\bottomrule
\end{longtable}

\section{Yêu cầu phi chức năng}

\begin{itemize}
  \item Giao diện nhất quán, tải lười các view lớn và giữ trạng thái form khi đổi tab.
  \item Tác vụ mạng dài chạy ở nền, có thông báo tiến trình và kết quả.
  \item Dev-mode không được mở session thật; nếu không xác minh được cờ dev thì phải fail-closed.
  \item Dữ liệu được kiểm tra định dạng, dùng khóa ngoại và giao dịch SQLite.
  \item Kiến trúc tách UI, bridge, persistence, worker và template.
  \item Password không được đưa vào log; mã hóa secret và phân quyền là yêu cầu tiếp theo.
\end{itemize}

\section{Kiến trúc tổng thể}

\begin{figure}[H]
\centering
\begin{tikzpicture}[
  node distance=7mm,
  box/.style={draw=ptitRed, rounded corners, align=center, minimum width=0.82\textwidth, minimum height=10mm, fill=red!3},
  arrow/.style={-{Latex[length=2.5mm]}, thick, draw=black!65}
]
  \node[box] (ui) {Qt Quick/QML UI\\Shell, Device, Interface, DHCP, Routing, ACL, NAT};
  \node[box, below=of ui] (bridge) {PyQt6 Bridge/Core\\DatabaseManager, runtime services, background task, View \& Push factory};
  \node[box, below=of bridge] (data) {Backend nghiệp vụ và SQLite\\normalize, validate, persistence, pending/applied state};
  \node[box, below=of data] (worker) {Network workers\\collector, Jinja2 renderer, connector, result updater};
  \node[box, below=of worker] (device) {Cisco IOS lab hoặc dev-mode};
  \draw[arrow] (ui) -- node[right]{signal/slot} (bridge);
  \draw[arrow] (bridge) -- (data);
  \draw[arrow] (data) -- (worker);
  \draw[arrow] (worker) -- (device);
\end{tikzpicture}
\caption{Kiến trúc runtime của \SoftwareName}
\label{fig:runtime-architecture}
\end{figure}

Entry point \pathcode{app/main.py} nạp module QML \codeinline{UI} và đưa các service Python vào context. \codeinline{DatabaseManager} là facade tổng hợp các mixin DHCP, ACL, NAT và các hàm thiết bị/routing.

\section{Luồng dữ liệu cốt lõi}

\subsection{Quản lý và đồng bộ thiết bị}

Người dùng tạo thiết bị, mở tab, ping hoặc kết nối. Connector lưu running-config theo host và parser đồng bộ một phần interface/OSPF vào database. Trạng thái tab quyết định có cho phép cấu hình hay không.

\subsection{Lưu cấu hình mong muốn}

Form QML kiểm tra dữ liệu cơ bản rồi gọi slot. Backend chuẩn hóa và ghi bản ghi với trạng thái pending. Thao tác xóa thường dùng soft delete để worker còn biết lệnh cần gỡ khỏi thiết bị.

\subsection{View và Push}

\begin{enumerate}
  \item Controller thu thập bản ghi pending theo host/module.
  \item Worker hoặc renderer tạo lệnh từ Jinja2 template.
  \item Hộp thoại hiển thị preview; chưa có lệnh thì khóa nút Push.
  \item Worker dùng connector của session hiện có hoặc xử lý dev-mode.
  \item Báo cáo thành công được dùng để chuyển pending thành applied hoặc xóa bản ghi chờ gỡ.
\end{enumerate}

\section{Thiết kế cơ sở dữ liệu}

\begin{table}[H]
\centering
\caption{Phân nhóm schema SQLite}
\label{tab:schema-groups}
\begin{tabularx}{\textwidth}{>{\bfseries}p{0.15\textwidth}X p{0.25\textwidth}}
\toprule
Tiền tố & Nội dung & Mức sử dụng \\
\midrule
\codeinline{t01} & Thiết bị và cấu hình truy cập/YANG & Runtime chính \\
\codeinline{t02} & Interface router L3, tunnel, WAN, QoS & UI/persistence, chưa push \\
\codeinline{t03} & DHCP pool, excluded, helper & Runtime chính \\
\codeinline{t04} & Static route, OSPF, EIGRP & Runtime chính, còn lệch test schema \\
\codeinline{t05} & ACL, binding, NAT, route-map & Runtime chính; ACL chưa push \\
\codeinline{t06} & Switching L2 & Schema dự phòng \\
\codeinline{t07} & VRF & Schema dự phòng \\
\bottomrule
\end{tabularx}
\end{table}

Trường \codeinline{success} thường mang ba trạng thái: 0 là pending thêm/sửa, 1 là đã áp dụng và $-1$ là pending remove/soft delete. Một số module còn dùng \codeinline{action_Cfg} để mô tả phần cấu hình thay đổi.

\section{Thiết kế giao diện}

UI gồm shell, activity bar, sidebar thiết bị, feature bar, content area và status bar. Các view lớn được tạo bằng Loader ở lần truy cập đầu và giữ sống để tránh mất dữ liệu form. Nhóm form/list được tái sử dụng cho DHCP/NAT, còn OSPF/EIGRP sử dụng process card và các section chuyên biệt.

\section{Thiết kế an toàn và xử lý lỗi}

Các worker kiểm tra dev-mode trước khi yêu cầu session thật. Tác vụ mạng có kết quả cấu trúc và chỉ cập nhật database khi worker báo thành công. Những nội dung chưa hoàn chỉnh gồm mã hóa credential, audit log, verify sau push và rollback giao dịch nhiều thiết bị.
